\section{Diskussion}\label{diskussion}

\textit{Die vollständige Diskussion wird ergänzt, sobald die experimentellen Ergebnisse vorliegen. Bereits jetzt belegbar sind die folgenden methodischen Bezugspunkte für die spätere Diskussion.}

\textit{Der Vergleich mit Raturi et al. \cite{raturi2026exploratory}, die die Balanced Accuracy als Metrik einführen, aber keine Erklärbarkeitskomponente integrieren. Der Vergleich mit Almahaqeri et al. \cite{almahaqeri2026gradient}, die auf demselben Datensatz mit Gradient Boosting arbeiten, jedoch stratifiziertes Undersampling statt einer hierarchischen Klassifikationsarchitektur wählen. Der Vergleich mit Alharby \cite{alharby2025ciciot}, der den PCA-Einsatz auf dem CICIoT2023-Datensatz dokumentiert und Vergleichswerte für die Ablation Study liefert. Die Bewertung der SHAP-Erklärungsqualität anhand der Metriken aus Hermosilla et al. \cite{hermosilla2025xai_forensic}, namentlich Fidelität, Konsistenz und Jaccard-Ähnlichkeit.}

\subsection*{Offene Forschungsfragen}

\textit{Dieser Abschnitt wird nach Vorliegen der Ergebnisse ausgeführt. Vorläufig identifizierte offene Fragen umfassen: die Latenzmessung und Echtzeit-Implementierung des Klassifikators, die Robustheit gegenüber gezielten Adversarial Attacks auf das ML-System sowie die Generalisierbarkeit der SHAP-Erklärungen auf weitere IoT-Datensätze.}
